\section{Marked periods and a complete local torsion criterion}\label{sec:periods}
The natural monodromy valuation is only one part of a toric period. The unit part can contain the global obstruction. Blocks 59 and 60 computed the marked period constants rather than replacing them by unspecified units.

At a variable place away from $2,3,5,7$, write $C=q^cC_0$ with $C_0$ a unit. In the marked two node basis, the actual leading period coordinates are
\begin{align*}
\Pi_{aa}&=-\frac{\eps\phi^2}{675\sqrt5}(1+\mathfrak m),\\
\Pi_{bb}&=\frac{\eps\phi^{-2}}{675\sqrt5}(1+\mathfrak m),\\
\Pi_{ab}&=\Pi_{ba}=-\phi^{-4}(1+\mathfrak m),
\end{align*}
where $\eps=C^7/A^3$ and the notation means that the quotient by the displayed leading term is a principal unit. The analytic uniformization and marked smoothing comparison are inherited written inputs, with precise sources in the literature report \cite{mumford1972,ichikawa1993}. The principal unit qualification is essential at seven, where those units may carry higher information. Away from seven they are divisible by every power of seven.

Put $M=7^j$ and write $[u]$ additively in $F^\times/F^{\times M}$ for the local field $F=K_v$. Let
\[
X=[\phi],\qquad Z=7[C]-3[A]-3[3]-5[\sqrt5].
\]
The entire marked period matrix becomes
\begin{equation}\label{eq:periodmatrix}
([\Pi_{ij}])=
\begin{pmatrix}Z+2X&-4X\\-4X&Z-2X\end{pmatrix}
=ZI+2XR,\qquad R=\begin{pmatrix}1&-2\\-2&-1\end{pmatrix},\quad R^2=5I.
\end{equation}
Since four is invertible modulo $M$, every marked period coordinate vanishes if and only if $X=Z=0$. The matrix $R$ belongs to this marked local period calculation. It is not silently identified with a global endomorphism or with an integral semilinear descent operator.

\begin{proposition}[Full torsion and period powers]\label{prop:torsionperiods}
For the stipulated split uniformization $J^{\mathrm{an}}=\mathbb G_m^2/\Lambda$ with rational period lattice, full $J[M]$ is $F$ rational if and only if $\mu_M\subset F$ and every coordinate of a basis of $\Lambda$ is an $M$th power in $F$.
\end{proposition}
\begin{proof}
The uniformization gives
\[
0\longrightarrow\mu_M^2\longrightarrow J[M]
\longrightarrow\Lambda/M\Lambda\longrightarrow0.
\]
If all torsion points are rational, so are $\mu_M^2$. For $\lambda\in\Lambda$, choose an $M$th root $u$ in the torus. Rationality of its torsion class gives $g(u)/u\in\Lambda$ for every $g\in G_F$. Since $\lambda$ is rational, this quotient also lies in $\mu_M^2$. The lattice is torsion free, so the intersection is trivial. Thus $u$ is rational. Every period coordinate is an $M$th power.

Conversely, take rational $M$th roots of each period basis vector. Their classes, together with $\mu_M^2$, generate all $M$ torsion by the exact sequence. Every generator is rational.
\end{proof}
At residue characteristic different from seven, a unit is an $M$th power precisely when its residue is, by simple root lifting. A general element also needs valuation divisible by $M$. If $Q$ is the residue cardinality, $\mu_M\subset F$ precisely when $M\mid Q-1$. Applying Proposition~\ref{prop:torsionperiods} to \eqref{eq:periodmatrix} gives the complete local test
\begin{equation}\label{eq:full-local-torsion}
\begin{gathered}
7^j\mid Q-1,\qquad7^j\mid7c,\\
\phi\in(\kappa_v^\times)^{7^j},\qquad
\frac{C_0^7}{3^3A^3(\sqrt5)^5}\in(\kappa_v^\times)^{7^j}.
\end{gathered}
\end{equation}
The valuation condition on $c$ is not lost when replacing $C$ by its unit part $C_0$.

At level seven, the last unit condition reduces to $75A$ being a seventh power. At level forty nine, exact multiplicative identities reduce it to
\[
3\cdot5^9A/C_0^{35}\in(\kappa_v^\times)^{49}.
\]
Thus full forty nine torsion is rational if and only if $49\mid Q-1$, $7\mid c$, $\phi$ is a forty ninth power residue, and the displayed rational unit is as well. These are two way local criteria for the stipulated local Jacobian.

They do not produce one global beta class or determine its actual depth. The archive retains positive local Frey families at $197,883,1373,1667$ satisfying the full criterion. Their survival is compatible with the later exclusion of a global beta supported only at those primes.

That exclusion illustrates why the common global coefficients matter. Correct inverse cyclotomic reciprocity on this four prime support leaves only inertia multiples of $(1,0,0,3)$. A nonzero class must therefore ramify at $197$ and $1667$. The natural rows in $(\phi,3,\sqrt5)$ coordinates are $(0,6,6)$ and $(0,3,4)$. They force both
\[
6a_3+6a_5=0,\qquad3a_3+4a_5=0.
\]
The determinant is six modulo seven, so $a_3=a_5=0$, contradicting \eqref{eq:globalalpha}. No value of $a_\phi$ can repair the contradiction because both unit symbols vanish. This is a global common coefficient exclusion. It does not retract any local positive family or prohibit an individual prime from appearing in a larger conductor.
